\subsection{Experimental setup}
\label{subsec:experiment-setup}
We assess both time discretizations by a manufactured-solution test
on $\Omega=(0,1)^3$ with $T=2$. Each Cartesian partition into
$K\times K\times K$ cubes is subdivided into tetrahedra, with
$K=4,8,16,32$ and $\dt=1/K=h/\sqrt{3}$, where $h$ is the maximum
tetrahedron diameter. Space and time are refined together;
this experiment does not isolate the temporal convergence order.

\paragraph{Manufactured solution.}
Define
\[
q(s)=64s^3(1-s)^3,\qquad
B(x,y,z)=q(x)q(y)q(z),\qquad
S(x,y,z)=\sin^2(\pi x)\sin^2(\pi y)\sin^2(\pi z),
\]
and let $a=(1,2,3)^\intercal$, $b=(1,-1,1/2)^\intercal$, and
$c=(1,2,-1)^\intercal$. The prescribed fields are
\begin{align*}
\bs u&=\frac{1+0.1t}{C_u}\,\nabla B\times a,
&\bs\omega&=\frac{1-0.05t}{C_\omega}\,B b,
&\bs m&=\frac{1+0.05t}{C_m}\,S c,\\
\varphi&=\frac{1-0.05t}{C_\varphi}\,B,
&\bs H&=\nabla\varphi,
&p&=\frac{1+0.05t}{C_p}\prod_{s\in\{x,y,z\}}\cos(2\pi s),
\end{align*}
where the fixed normalization constants used in the computations are
\begin{gather*}
C_u=363.275188821,\quad C_\omega=20.6500992827,\quad
C_m=\frac{\sqrt{81+108\pi^2}}{16},\\
C_\varphi=13.7667328551,\qquad C_p=0.353553390593.
\end{gather*}
The auxiliary fields are $\bs k=\curl\bs m$, $\bs z=\bs u\times\bs m$,
$\vartheta=-\Delta\varphi$, and $\varpi=\bs u\cdot\bs m$.
We set $\bs H_e=-\mu_0(\bs m+\bs H)$ and use the zero-mean modified
pressure $\widetilde p=\mathcal P_0[p-\mu_0\bs m\cdot\bs H/2]$, where
$\mathcal P_0v=v-|\Omega|^{-1}\int_\Omega v$.
All physical parameters $\rho,\kappa,\eta,\zeta,\mu_0,\sigma,
\eta^\prime,\lambda^\prime,\tau,\chi_0$ are set to one.

This is a forced verification problem. The momentum, angular-momentum,
and magnetization sources $\bs f,\bs g,\bs r$ are obtained by substituting
the prescribed fields into the corresponding continuous equations.
Their weak right-hand sides are $(\bs f,\bs v_h)$,
$(\bs g,\bs s_h)$, and $(\bs r,\bs F_h)$; the magnetic-potential
right-hand side is
\[
\left(-\partial_t\bs H_e
-\frac{\mu_0+\chi_0}{\tau\mu_0}\bs H_e
+\sigma\nabla\div\bs H_e-\mu_0\bs r,\nabla\psi_h\right).
\]
The fields satisfy the homogeneous boundary traces used in the experiment:
zero velocity and angular velocity, zero normal magnetization,
zero tangential $\bs z,\bs k,\bs H$, and zero
$\varphi,\vartheta,\varpi$. In particular, $q,q',q''$ vanish at both
endpoints and $S,\nabla S$ vanish on the cube boundary. The new
$\vartheta$ and $\varpi$ couplings are not identically zero for this test.

\paragraph{Discretization and implementation.}
The first-order runs use the MINI velocity element, continuous piecewise
linear pressure and angular velocity, lowest-order Raviart--Thomas
magnetization and first-family Nedelec $\bs z,\bs k,\bs H$, and
continuous piecewise linear $\varphi,\vartheta,\varpi$.

For the second-order scheme, velocity and angular velocity use continuous
piecewise quadratic elements, pressure remains continuous piecewise linear, magnetization
uses RT1, and $\bs z,\bs k,\bs H$ use first-family Nedelec elements
one degree above the lowest order. The scalar fields
$\varphi,\vartheta,\varpi$ are continuous piecewise quadratic.
The corresponding RT and Nedelec Basix degree is two.

The initial velocity is obtained by a discrete Stokes projection.
After interpolating the initial magnetization into its RT space,
the initial potential is determined from the compatible magnetic Poisson
problem. The initial $\bs k,\bs z,\vartheta,\varpi$ satisfy their discrete
mass-projection relations. Pressure is normalized by its finite-element
integral mean.

For the second-order convergence test, the first time step is solved
by a Broyden quasi-Newton iteration with
relative tolerance $10^{-9}$ and a maximum of 12 iterations.
The accepted state is checked by reassembling the nonlinear midpoint
equations. Subsequent time steps use the prescribed extrapolated
coefficients and require one five-block linear solve per step.
The $\bs z_h$ and $\varpi_h$ slots represent the current midpoint-stage
projections; they are not averaged again with the preceding stage.

Computations use FEniCSx/DOLFINx~0.10.0~\cite{Barattafenics}
with PETSc and MUMPS. The five-block Algorithm~3 is evaluated through
matrix-free Schur-complement actions. The $K_1,K_3,K_5$ subproblems use
MUMPS LU; the $S_2,S_4$ systems use FGMRES with LU preconditioners.
The MUMPS ordering control is $\mathrm{ICNTL}(7)=7$.
The Schur tolerances are $10^{-10}$ relative and $10^{-12}$ absolute.
All first-order cases and the second-order cases up to $K=16$ use
48 MPI processes with one thread each; the second-order $K=32$ case
uses 32 MPI processes with four threads each on a 1-TB-class node.
Source terms are evaluated at quadrature points, with assembly
quadrature degree 12. In the error calculations, the analytic values
and derivatives are independently represented in a degree-6 DG space
and integrated with degree-12 quadrature; the numerical solution is
evaluated in its original finite-element space.

As implementation checks for the manufactured-solution runs,
the two schemes were compared with a
monolithic solve of the same matrix at the first two time levels on
$K=4$, modulo the same pressure
nullspace; the normalized solution differences were below $1.1\times10^{-12}$.
An independent transcription of the nine weak equations on $K=2$,
using random finite-element fields and both unit and non-unit parameters,
also agreed with the production matrices and right-hand sides to
scaled differences below $1.1\times10^{-16}$.
These checks concern algebraic and implementation consistency, rather
than replacing the convergence test.

\paragraph{Error measures.}
For a field $v$ and its specified norm $X$, we report
\[
E_X(v)=\frac{\|v-v_h\|_X}{\|v\|_X},\qquad
\mathrm{rate}(K\to2K)=\log_2\frac{E_X(v;K)}{E_X(v;2K)}.
\]
The norms are full $H^1$ for $\bs u,\bs\omega,\varphi$, $L^2$ for
$\widetilde p,\bs k$, $H(\div)$ for $\bs m$, and $H(\curl)$ for
$\bs H$. Thus the graph norms combine the squared $L^2$ and derivative
terms before taking the square root. Pressure errors use zero-mean
numerical and exact fields. First-order errors are evaluated at $T$.

For the second-order scheme, the primary fields are compared at $T$.
Its pressure slot represents the averaged modified pressure, so its reference is
\[
\widetilde p_{\mathrm{ref}}^{N-1/2}
=\frac{\widetilde p(T)+\widetilde p(T-\dt)}{2},
\]
with the same zero-mean convention as above. This reference is not
identified with either the endpoint pressure or the analytic pressure
evaluated directly at the midpoint.

We also report $C_H$, the maximum of $|\curl\bs H_h|$ over the
degree-6 DG interpolation points, as a sampled estimate of
$\|\curl\bs H_h\|_{L^\infty}$. It is an absolute diagnostic, not a
relative error or a rigorous continuous supremum.
